\section{Вывод}

Мы нашли радиус кривизны линзы через квадрат диаметров колец Ньютона, оценили погрешность вычисления. Качественно исследовали влияние углового размера источника света на интерференционную картину - это обратная зависимость (больше угловой размер источника - меньше наблюдаемых колец). Пронаблюдали кольца Ньютона в белом свете. Они представляют собой разноцветные кольца, идущие в радужном порядке, сливающиеся в белый цвет, при отдалении от центра.